%Plantilla basada en "Template for Masters / Doctoral Thesis" (plantilla disponible en writeLaTex) que subió LaTeXTemplates.com

\documentclass[11pt]{book}
\usepackage[paperwidth=17cm, paperheight=22.5cm, bottom=2.5cm, right=2.5cm]{geometry}
\usepackage{amssymb,amsmath,amsthm} %paquete para símbolo matemáticos
\usepackage[spanish]{babel}
\usepackage[utf8]{inputenc} %Paquete para escribir acentos y otros símbolos directamente
\usepackage{enumerate}
%\usepackage{graphicx}
\usepackage[pdftex]{color,graphicx}   
\usepackage{tikz}
\usepackage{float}
\usepackage{url}
%\usepackage[]{algorithm2e}
\usepackage[linesnumbered,ruled,vlined]{algorithm2e}


\newenvironment{Algoritmo}[1][htb]
  {\renewcommand{\algorithmcfname}{Algoritmo}% Update algorithm name
   \begin{algorithm}[H]%
  }{\end{algorithm}}
  
 \usepackage{algcompatible}

\usetikzlibrary{arrows}
\usetikzlibrary{quotes}
\usetikzlibrary{shapes,snakes}
\tikzset{cross/.style={cross out, draw=black, minimum size=2*(#1-\pgflinewidth), inner sep=0pt, outer sep=0pt},
%default radius will be 1pt. 
cross/.default={10pt}}
\newcommand{\tcancel}[2][black]{%
\begin{tikzpicture}
\node[draw=#1,cross out,inner sep=1pt] (a){#2};
\end{tikzpicture}%
}


% Gian Carlo Diluvi's preamble ---------
\usepackage{amsthm}
\usepackage{amsmath}
\usepackage[nottoc]{tocbibind}
\usepackage{apacite} % Citar con APA
%\usepackage[numbers]{natbib}
\usepackage{natbib}
%\usepackage{biblatex}
%\bibliography{bibliografia}
\usepackage{here}
\usepackage{float}
\usepackage{etoolbox}
\usepackage{url}
\usepackage{mathrsfs}
\usepackage{commath}
\usepackage{color}
\usepackage{multirow}
%\usepackage{subcaption}



%\usepackage{subfig} %para poner subfiguras
\graphicspath{{Img/}} %En qué carpeta están las imágenes
\usepackage[nottoc]{tocbibind}

\usepackage{listings}
	\lstset{frame = single,
    		literate={á}{{\'a}}1
        			 {ã}{{\~a}}1
        			 {é}{{\'e}}1
        			 {ó}{{\'o}}1
        			 {í}{{\'i}}1
        			 {ñ}{{\~n}}1
        			 {¡}{{!`}}1
        			 {¿}{{?`}}1
        			 {ú}{{\'u}}1
                     {Á}{{\'A}}1
                     {É}{{\'E}}1
                     {Í}{{\'I}}1
                     {Ó}{{\'O}}1
                     {Ú}{{\'U}}1
}
\renewcommand{\lstlistingname}{C\'odigo}
\lstdefinestyle{R}{%
%\captionsetup{labelformat=algocaption,labelsep=colon}
    mathescape=false,
    breaklines=true,
    frame=single,
    numbers=left, 
    numberstyle=\small,
    language=R,
    basicstyle=\scriptsize\ttfamily, 
    keywordstyle=\color{black}\bf,
    xleftmargin=.04\textwidth,
}

\usepackage[pdftex,
            pdfauthor={Ilan Jinich Fainsod},
            pdftitle={El problema de admisión a universidades con cotas inferiores},
            pdfsubject={ÁREA DE LA TESIS},
            pdfkeywords={PALABRAS CLAVE},
            pdfproducer={Latex con hyperref},
            pdfcreator={pdflatex}]{hyperref}



\begin{document}

%----------------------------------------------------------------------------------------
%	COMANDOS PERSONALIZADOS
%----------------------------------------------------------------------------------------

%SI TU TESIS TIENE TEOREMAS Y DEMOSTRACIONES, PUEDES DESCOMENTAR Y USAR LOS SIGUIENTES COMANDOS

\renewcommand{\proofname}{Demostración}
%\providecommand{\norm}[1]{\lVert#1\rVert} %Provee el comando para producir una norma.
%\providecommand{\innp}[1]{\langle#1\rangle} 
%\newcommand{\seno}{\mathrm{sen}}
%\newcommand{\diff}{\mathrm{d}}

\newtheorem{teo}{Teorema}[section] 
\newtheorem{cor}[teo]{Corolario}
\newtheorem{lem}[teo]{Lema}
%\renewcommand{\qedsymbol}{$\blacksquare$}
%\renewcommand{\qedsymbol}{$\dag$}}
\newcommand{\fin}{\hfill$\triangle$}

\theoremstyle{definition}
\newtheorem{dfn}[teo]{Definición}

\theoremstyle{remark}
\newtheorem{obs}[teo]{Observación}

\newtheorem{eje}{Ejemplo}[section]

\allowdisplaybreaks


%----------------------------------------------------------------------------------------
%	PORTADA
%----------------------------------------------------------------------------------------

\title{TÍTULO DE LA TESIS} %Con este nombre se guardará el proyecto en writeLaTex

\begin{titlepage}
\begin{center}

\textsc{\Large Instituto Tecnológico Autónomo de México}\\[4em]

%Figura
\begin{figure}[h]
\begin{center}
\includegraphics{logo-ITAM_ch.jpg}
\end{center}
\end{figure}

\vspace{4em}

\textsc{\huge \textbf{El problema de admisión a universidades con cotas inferiores y comunes}}\\[4em]

\textsc{\large Tesis}\\[1em]

\textsc{que para obtener el título de}\\[1em]

\textsc{LICENCIADO EN MATEMÁTICAS APLICADAS}\\[1em]

\textsc{presenta}\\[1em]

\textsc{\Large Ilan Jinich Fainsod}\\[1em]

\textsc{\large Asesor: DR RAMÓN ESPINOSA ARMENTA }

\end{center}

\vspace*{\fill}
\textsc{México, D.F. \hspace*{\fill} 2018}

\end{titlepage}


%----------------------------------------------------------------------------------------
%	DECLARACIÓN
%----------------------------------------------------------------------------------------

\thispagestyle{empty}
\vspace*{\fill}
\begingroup
``Con fundamento en los artículos 21 y 27 de la Ley Federal del Derecho de Autor y como titular de los derechos moral y patrimonial de la obra titulada ``\textbf{El problema de admisión a universidades con cotas inferiores y comunes}'', otorgo de manera gratuita y permanente al Instituto Tecnológico Autónomo de México y a la Biblioteca Raúl Bailléres Jr., la autorización para que fijen la obra en cualquier medio, incluido el electrónico, y la divulguen entre sus usuarios, profesores, estudiantes o terceras personas, sin que pueda percibir por tal divulgación una contraprestación''.

\centering

\hspace{3em}

\textsc{AUTOR}

\vspace{5em}

\rule[1em]{20em}{0.5pt} % Línea para la fecha

\textsc{Fecha}
 
\vspace{8em}

\rule[1em]{20em}{0.5pt} % Línea para la firma

\textsc{Firma}

\endgroup
\vspace*{\fill}


%----------------------------------------------------------------------------------------
%	DEDICATORIA
%----------------------------------------------------------------------------------------

\pagestyle{empty}
\frontmatter

\chapter*{}
\begin{flushright}
%\textit{``To life, to life, L'Chaim!''}
\textit{$6\times 10^6$}
%\textit{זאָג ניט קיין מאָל}
\end{flushright}


%----------------------------------------------------------------------------------------
%	AGRADECIMIENTOS
%----------------------------------------------------------------------------------------

\chapter*{Agradecimientos}
%\markboth{AGRADECIMIENTOS23}{AGRADECIMIENTOS} % encabezado 

¡Muchas gracias a todos!


%----------------------------------------------------------------------------------------
%	PREFACIO
%----------------------------------------------------------------------------------------

\chapter*{Prefacio}

\pagestyle{plain}
%\markboth{PREFACIO23}{PREFACIO} % encabezado 

PUEDEN QUITAR ESTA PARTE


%----------------------------------------------------------------------------------------
%	TABLA DE CONTENIDOS
%---------------------------------------------------------------------------------------

\tableofcontents
\newpage


%----------------------------------------------------------------------------------------
%	TESIS
%----------------------------------------------------------------------------------------
\mainmatter %empieza la numeración de las páginas
\pagestyle{headings}

%  Incluye los capítulos en el folder de capítulos

\chapter{Introducción}
\begin{flushright}
\textit{``Matchmaker, matchmaker,
Make me a match, Find me a find,
Catch me a catch''}
\end{flushright}
\section{Descripción del problema}
Imaginemos que en una ciudad tenemos $m$ universidades y $n$ personas que desean ir a una de estas escuelas, cada universidad tiene un límite alumnos que puede admitir y por esto no todas las personas pueden asistir a su universidad ''favorita''. Además, podemos agregar el supuesto de que las escuelas tienen cotas inferiores, usualmente por razones de dinero o de personal a una escuela no le conviene abrir o aceptar a menos de que entren una cantidad considerable de alumnos a esta, este tipo de cota puede ser común si dos o más escuelas comparten recursos o personal. Adicionalmente, tenemos el supuesto de que todos los alumnos tienen una lista de preferencias exhaustiva de que universidad prefieren y de forma análoga las universidades tienen una lista de preferencias sobre los alumnos, cada escuela y persona tienen necesidades diferentes a los demás, algunos podrían preferir la escuela que está más cerca de sus casas, otros la que te prepara mejor académicamente y otros entre muchas razones podrían preferir la que tenga el mejor ambiente social. 
\\ En notación matemática lo que hacemos es suponer que tenemos $n$ alumnos $\alpha_1, \alpha_2, \ldots, \alpha_n$ y $m$ universidades $a_1, a_2, \ldots, a_m$, con las restricciones 

\begin{enumerate}
\item \begin{equation} \label{r1}
x_{i,j}= 
\begin{cases}
1 & \qquad \text{si $i$ asiste a la universidad $j$.} \\
0 &\qquad\text{en otro caso.}\ \\ 
\end{cases} \end{equation}
\item \begin{equation} y_{j}= 
\begin{cases}
1 & \qquad \text{si la universidad $j$ abre.} \\
0 &\qquad\text{en otro caso.} \\ 
\end{cases} \end{equation}
\item \begin{equation} \label{r2}
\sum_{j=1}^{m}x_{i,j} \leq1 \ \text{ para toda $i=1,2,\ldots,n$. }
\end{equation} Cada estudiante es admitido en máximo una universidad. 
\item \begin{equation} \label{r3}
\sum_{i=1}^{n} x_{i,j} \leq M_j\ \text{ para toda $j=1,2,\dots,m$.} 
\end{equation}
Cada universidad tiene un límite de alumnos que puede admitir.
\item \begin{equation} \label{r4}
\sum_{i=1}^{n} x_{i,j} \geq\ m_j\times y_j 
\end{equation}
Cada universidad necesita admitir a una cantidad considerable de alumnos para abrir.
\item \begin{equation} \label{r5}
\sum_{i=1}^{n} (x_{i,j}+x_{i,j'}) 
\begin{cases}
\geq m_{j,j'} \times y_j \\
\geq m_{j,j'} \times y_{j'} \\
\end{cases}
\end{equation} \\
para toda $(j,j')$, con $j$ distinto de $j'$. \\
La cota inferior de las universidades puede ser común.
\item \begin{equation} \label{r6}
x_{i,j} \leq y_j \ \text{ para toda $i=1,2,\ldots,n$ y para toda $j=1,2,\ldots,m$.}
\end{equation}
Los solicitantes únicamente pueden asistir a universidades abiertas.

\end{enumerate}

Es importante mencionar que no hay garantía de que todos los alumnos sean admitidos en alguna universidad, algunos se pueden quedar sin estudiar. No todas las universidades necesariamente tienen cotas superiores, inferiores o comunes; en ciertas situaciones se puede considerar que para cierta universidad $j$, $M_{j}$ puede ser igual a infinito o $m_j$ puede ser igual a cero y de igual manera para cierto par de universidades $(j,j')$, $m_{j,j'}$ puede ser igual a cero. 

Para definir las preferencias de una forma un poco más formal, hacemos uso de la matriz de preferencias

\begin{dfn}
\label{matpref}
Definimos la \textbf{matriz de preferencias} para un problema con $n$ estudiantes y con $m$ universidades como una matriz con $n$ filas y $m$ columnas donde la primera entrada de la celda $(i,j)$ represente el orden de preferencia que le asigna el solicitante $i$ a la universidad $j$ y análogamente la segunda entrada representa el orden de preferencia que le da la universidad $j$ a la persona $i$.
\end{dfn}

Para que dejar más clara la definición es necesario mostrar un ejemplo. 

\begin{eje}
\label{ejemplo matrimonio 1}
Supongamos que contamos con tres solicitantes $\alpha$, $\beta$ y $\gamma$ y con tres universidades $A$, $B$ y $C$ y la matriz de preferencias
$$\begin{pmatrix}
& A & B & C \\
\alpha & 1,3 & 2,2 & 3,1 \\
\beta & 3,1 & 1,3 & 2,2 \\
\gamma & 2,2 & 3,1 & 1,3 
\end{pmatrix}$$
aquí por ejemplo el orden preferencias de $\alpha$ es $(A,B,C)$ y el orden de preferencias de $C$ es $(\gamma, \beta, \alpha)$.
\fin
\end{eje}


Para continuar es necesario introducir un criterio de estabilidad y otro de optimalidad, para ello enunciamos dos definiciones. 

\begin{dfn}{\cite{GaleShapley}}
\label{Estable}
Decimos que una asignación de solicitantes a universidades es inestable si existen dos solicitantes $\alpha$ y $\beta$ asignados a universidades $A$ y $B$ respectivamente, con la propiedad de que $\alpha$ prefiere estar en $B$ que en $A$ y $B$ prefiere tener a $\alpha$ que a $B$. Es decir, existen una universidad y un solicitante que se prefieren entre ellos a sus respectivas asignaciones. Alternativamente una asignación de solicitantes a universidades es \textbf{estable} si no es inestable.
\end{dfn}

\begin{dfn}{\cite{GaleShapley}}
\label{optima}
Una asignación es considera \textbf{optima} si cada solicitante esta mejor o igual respecto a sus preferencias que en cualquier otra asignación estable. 
\end{dfn}

Un par de ejemplos de estos conceptos son 

\begin{eje}
Retomando el ejemplo \ref{ejemplo matrimonio 1}. \\
El conjunto de parejas $(\alpha, A), \; (\beta, B)\; y\; (\gamma, C)$ es estable porque a pesar de las preferencias de las universidades cada persona esta con su primera opción, es decir, los solicitantes prefieren a su universidad más que a cualquier otra. Esta además es óptima porque cada solicitante está en su primera opción.
\begin{figure}[H]\centering

\begin{tikzpicture}[ scale=0.8]
\tikzset{vertex/.style = {shape=circle,draw,minimum size=1.5em}}
\tikzset{edge/.style = {->,> = latex}}
\filldraw[color=blue!60, fill=blue!5, very thick](0,3) ellipse (.8 and 1.5);
\filldraw[color=red!60, fill=red!5, very thick](4,3) ellipse (.8 and 1.5);


% vertices
% 


\node[vertex] (a) at (0,4) {$\alpha$};
\node[vertex] (b) at (0,3) {$\beta$};
\node[vertex] (c) at (0,2) {$\gamma$};


\node[vertex] (e) at (4,4) {$A$};
\node[vertex] (f) at (4,3) {$B$};
\node[vertex] (g) at (4,2) {$C$};


\node (i) at (0,5) {Hombres};
\node (j) at (4,5) {Mujeres};

\path[-stealth] (a) edge (e);
\path[-stealth] (b) edge (f);
\path[-stealth] (c) edge (g);



%\draw (0.2,8)--(3.8,8);



\end{tikzpicture}

\caption{Emparejamiento estable.}
\end{figure}

\fin
\end{eje}
\begin{eje}
Retomando el ejemplo \ref{ejemplo matrimonio 1}. \\
El conjunto de parejas $(\alpha, A)$, $(\gamma, B)$ y $(\beta, C)$ es inestable porque $\gamma$ prefiere a $A$ que a su universidad actual ($B$) y A prefiere a $\gamma$ que a su universidad actual $(\alpha)$.

\begin{figure}[H]\centering

\begin{tikzpicture}[ scale=0.8]
\tikzset{vertex/.style = {shape=circle,draw,minimum size=1.5em}}
\tikzset{edge/.style = {->,> = latex}}
\filldraw[color=blue!60, fill=blue!5, very thick](0,3) ellipse (.8 and 1.5);
\filldraw[color=red!60, fill=red!5, very thick](4,3) ellipse (.8 and 1.5);


% vertices
% 


\node[vertex] (a) at (0,4) {$\alpha$};
\node[vertex] (b) at (0,3) {$\beta$};
\node[vertex] (c) at (0,2) {$\gamma$};


\node[vertex] (e) at (4,4) {$A$};
\node[vertex] (f) at (4,3) {$B$};
\node[vertex] (g) at (4,2) {$C$};


\node (i) at (0,5) {Hombres};
\node (j) at (4,5) {Mujeres};

\path[-stealth] (a) edge (e);
\path[-stealth] (c) edge (f);
\path[-stealth] (b) edge (g);



%\draw (0.2,8)--(3.8,8);



\end{tikzpicture}

\caption{Emparejamiento inestable.}
\end{figure}
\fin
\end{eje}


Surgen algunas preguntas de aquí ¿siempre existe un emparejamiento estable? ¿cómo se encuentra? Para resolverlas primero consideraremos el caso en el que $M_j=1$ para toda universidad $j$ y donde además no hay cotas inferiores o comunes, que es mejor conocido como el problema del matrimonio estable. 


%Descripción
%Definiciones
%Chane esto jala mejor como capitulo y no sección
%Comentario que explique cómo abajo es caso particular

\section{El problema del matrimonio estable}
%\begin{flushright}
%\textit{``Even the worst husband, God forbid, is better than no husband, God forbid.''}
%\end{flushright}
En el pequeño Shtetl de Anatevka, había muchos hombres y muchas mujeres jóvenes con el deseo de casarse. Yenta, la casamentera del pueblo, tenía el trabajo y la obligación de emparejar a todos estos jóvenes. Esta tarea sonara fácil pero no lo es, cada joven tenía distintas preferencias y cada uno de ellos necesitaba un trato especial, "el matrimonio es para toda la vida y no es algo que hay que tomar como si fuera cualquier cosa" decía Yenta. Retomando los conceptos de la última sección, se desea que el emparejamiento de estos jóvenes sea estable y optimo si es que es posible.\footnote{Este problema hace el supuesto de que las relaciones solo pueden ser entre un hombre y una mujer y que esto es reflejado por sus preferencias. Esto no manifiesta la realidad o las opiniones del autor y se hace unicamente con fines matemáticos.}
En terminos matemáticos, supongamos que tenemos $n$ hombres $$\alpha_1,\alpha_2,\ldots,\alpha_n$$ y $m$ mujeres $$a_1, a_2,\ldots,a_m$$ con las restricciones
\begin{enumerate}
\item \begin{equation} \label{1r1}
x_{i,j}= 
\begin{cases}
1 & \qquad \text{si $i$ se casa con $j$.} \\
0 &\qquad\text{en otro caso.}\ \\ 
\end{cases} \end{equation}
\item \begin{equation} \label{1r2}
\sum_{j=1}^{m}x_{i,j} \leq1 \ \text{ para toda $i=1,2,\ldots,n$. }
\end{equation} Cada hombre se casa con máximo una mujer. 
\item \begin{equation} \label{1r3}
\sum_{i=1}^{n} x_{i,j} \leq 1\ \text{ para toda $j=1,2,\dots,m$.} 
\end{equation}
Cada mujer se casa con máximo una hombre. 
\end{enumerate}

Es claro, a partir de la definición del problema que este es un caso particular del problema original. Aprovechando lo que se hizo en la sección anterior la definición de la matrz de preferencias, un emparejamiento estable y de un emparejamiento optimo son analogas a las definiciones \ref{matpref}, \ref{Estable} y \ref{optima}. Existen diversos resultados sobre el problema del matrimonio estable, uno de los más famosos es por Gale y Shapley que encontraron un algoritmo que siempre encuentra un emparejamiento estable y además lo hace sin muchas complicaciones. 

Un supuesto adicional que haremos es que el número de hombres es exactamente igual al numero de mujeres. Esto se vera despues que es relativamente sencillo de generalizar al caso en el que el número de hombres y de mujeres no es igual. 

%Definición del problema
%Definiciones alternativas
%

%Ejemplos

\subsection{Algoritmo de Gale Shapley}
En 1962 David Gale y Lloyd Stowell Shapley demostraron que para el problema del matrimonio estable siempre se puede encontrar un emparejamiento estable óptimo para los hombres, esto se hizo enunciando un algoritmo y mostrando su convergencia. El algoritmo queda representado por el siguiente codigo y es conocido como el algoritmo de Gale Shapley.

\begin{lstlisting}[style=R, escapeinside={(*}{*)},caption={Algoritmo de Gale Shapley}, captionpos=b, label=c1]
Input : Una matriz de preferencias para (*$n$*) hombres y (*$n$*) mujeres 
Output: Un emparejamiento. 
Cada hombre le propone a la primera mujer de su lista. Cada mujer que recibe más de una propuesta acepta la que este más arriba en su lista y rechaza al resto. 
repeat{ #hasta que todos los hombres tengan pareja
Los hombres no emparejados le proponen a la siguiente mujer en su lista. 
Cada mujer que recibe una propuesta escoge la que está más arriba en su lista entre las propuestas que recibió y su pareja actual.
Las mujeres rechazan las propuestas que no aceptaron.
}
\end{lstlisting}
Para dejar la explicación un poco más clara es importante mostrar un ejemplo de cómo funciona el algoritmo y cuál es su resultado final.
\begin{eje}{\cite{GaleShapley}}
\label{ejemploGS}
Supongamos que para 4 hombres y 4 mujeres tenemos la matriz de preferencias 
$$\begin{pmatrix}
& A & B & C & D \\
\alpha & 1,3 & 2,2 & 3,1 & 4,3 \\
\beta & 1,4 & 2,3 & 3,2 & 4,4 \\
\gamma & 3,1 & 1,4 & 2,3 & 4,2 \\
\delta & 2,2&3,1 &1,4 & 4,1 
\end{pmatrix}.$$
En primera estancia, $\alpha$ le propone matrimonio a $A$, $\beta$ le propone matrimonio a $A$, $\gamma$ le propone matrimonio a $B$ y $\delta$ le propone matrimonio a $C$.

\begin{figure}[H]\centering

\begin{tikzpicture}[ scale=0.8]
\tikzset{vertex/.style = {shape=circle,draw,minimum size=1.5em}}
\tikzset{edge/.style = {->,> = latex}}
\filldraw[color=blue!60, fill=blue!5, very thick](0,2.5) ellipse (.8 and 2);
\filldraw[color=red!60, fill=red!5, very thick](4,2.5) ellipse (.8 and 2);


% vertices
% 


\node[vertex] (a) at (0,4) {$\alpha$};
\node[vertex] (b) at (0,3) {$\beta$};
\node[vertex] (c) at (0,2) {$\gamma$};
\node[vertex] (d) at (0,1) {$\delta$};

\node[vertex] (e) at (4,4) {$A$};
\node[vertex] (f) at (4,3) {$B$};
\node[vertex] (g) at (4,2) {$C$};
\node [vertex] (h) at (4,1) {$D$};

\node (i) at (0,5) {Hombres};
\node (j) at (4,5) {Mujeres};

\path[-stealth] (a) edge (e);
\path[-stealth] (b) edge (e);
\path[-stealth] (c) edge (f);
\path[-stealth] (d) edge (g);



%\draw (0.2,8)--(3.8,8);



\end{tikzpicture}

\caption{Primera iteración.}
\end{figure}

En segunda instancia, $A$ acepta la propuesta de $\alpha$ y $\beta$ le propone matrimonio a $B$.

\begin{figure}[H]
\centering

\begin{tikzpicture}[ scale=0.8]
\tikzset{vertex/.style = {shape=circle,draw,minimum size=1.5em}}
\tikzset{edge/.style = {->,> = latex}}
\filldraw[color=blue!60, fill=blue!5, very thick](0,2.5) ellipse (.8 and 2);
\filldraw[color=red!60, fill=red!5, very thick](4,2.5) ellipse (.8 and 2);


% vertices
% 


\node[vertex] (a) at (0,4) {$\alpha$};
\node[vertex] (b) at (0,3) {$\beta$};
\node[vertex] (c) at (0,2) {$\gamma$};
\node[vertex] (d) at (0,1) {$\delta$};

\node[vertex] (e) at (4,4) {$A$};
\node[vertex] (f) at (4,3) {$B$};
\node[vertex] (g) at (4,2) {$C$};
\node [vertex] (h) at (4,1) {$D$};

%\draw (0,.5) node[cross,red] {};

\node (i) at (0,5) {Hombres};
\node (j) at (4,5) {Mujeres};

\path[-stealth] (a) edge (e);
\path[-stealth] (b) edge (f);
\path[-stealth] (c) edge (f);
\path[-stealth] (d) edge (g);



%\draw (0.2,8)--(3.8,8);



\end{tikzpicture}

\caption{Segunda iteración.}
\end{figure}


En tercera instancia, $B$ acepta la propuesta de $\beta$ y $\gamma$ le propone matrimonio a $C$.

\begin{figure}[H]\centering

\begin{tikzpicture}[ scale=0.8]
\tikzset{vertex/.style = {shape=circle,draw,minimum size=1.5em}}
\tikzset{edge/.style = {->,> = latex}}
\filldraw[color=blue!60, fill=blue!5, very thick](0,2.5) ellipse (.8 and 2);
\filldraw[color=red!60, fill=red!5, very thick](4,2.5) ellipse (.8 and 2);


% vertices
% 


\node[vertex] (a) at (0,4) {$\alpha$};
\node[vertex] (b) at (0,3) {$\beta$};
\node[vertex] (c) at (0,2) {$\gamma$};
\node[vertex] (d) at (0,1) {$\delta$};

\node[vertex] (e) at (4,4) {$A$};
\node[vertex] (f) at (4,3) {$B$};
\node[vertex] (g) at (4,2) {$C$};
\node [vertex] (h) at (4,1) {$D$};

\node (i) at (0,5) {Hombres};
\node (j) at (4,5) {Mujeres};

\path[-stealth] (a) edge (e);
\path[-stealth] (b) edge (f);
\path[-stealth] (c) edge (g);
\path[-stealth] (d) edge (g);



%\draw (0.2,8)--(3.8,8);



\end{tikzpicture}

\caption{Tercera iteración.}
\end{figure}

En cuarta instancia, $C$ acepta la propuesta de $\gamma$ y $\beta$ le propone matrimonio a $A$.

\begin{figure}[H]\centering

\begin{tikzpicture}[ scale=0.8]
\tikzset{vertex/.style = {shape=circle,draw,minimum size=1.5em}}
\tikzset{edge/.style = {->,> = latex}}
\filldraw[color=blue!60, fill=blue!5, very thick](0,2.5) ellipse (.8 and 2);
\filldraw[color=red!60, fill=red!5, very thick](4,2.5) ellipse (.8 and 2);


% vertices
% 


\node[vertex] (a) at (0,4) {$\alpha$};
\node[vertex] (b) at (0,3) {$\beta$};
\node[vertex] (c) at (0,2) {$\gamma$};
\node[vertex] (d) at (0,1) {$\delta$};

\node[vertex] (e) at (4,4) {$A$};
\node[vertex] (f) at (4,3) {$B$};
\node[vertex] (g) at (4,2) {$C$};
\node [vertex] (h) at (4,1) {$D$};

\node (i) at (0,5) {Hombres};
\node (j) at (4,5) {Mujeres};

\path[-stealth] (a) edge (e);
\path[-stealth] (b) edge (f);
\path[-stealth] (c) edge (g);
\path[-stealth] (d) edge (e);



%\draw (0.2,8)--(3.8,8);



\end{tikzpicture}

\caption{Cuarta iteración.}
\end{figure}


En quinta instancia, $A$ acepta la propuesta de $\delta$ y $\alpha$ le propone matrimonio a $B$.

\begin{figure}[H]\centering

\begin{tikzpicture}[ scale=0.8]
\tikzset{vertex/.style = {shape=circle,draw,minimum size=1.5em}}
\tikzset{edge/.style = {->,> = latex}}
\filldraw[color=blue!60, fill=blue!5, very thick](0,2.5) ellipse (.8 and 2);
\filldraw[color=red!60, fill=red!5, very thick](4,2.5) ellipse (.8 and 2);


% vertices
% 


\node[vertex] (a) at (0,4) {$\alpha$};
\node[vertex] (b) at (0,3) {$\beta$};
\node[vertex] (c) at (0,2) {$\gamma$};
\node[vertex] (d) at (0,1) {$\delta$};

\node[vertex] (e) at (4,4) {$A$};
\node[vertex] (f) at (4,3) {$B$};
\node[vertex] (g) at (4,2) {$C$};
\node [vertex] (h) at (4,1) {$D$};

\node (i) at (0,5) {Hombres};
\node (j) at (4,5) {Mujeres};

\path[-stealth] (a) edge (f);
\path[-stealth] (b) edge (f);
\path[-stealth] (c) edge (g);
\path[-stealth] (d) edge (e);



%\draw (0.2,8)--(3.8,8);



\end{tikzpicture}

\caption{Quinta iteración.}
\end{figure}

En sexta instancia, $B$ acepta la propuesta de $\alpha$ y $\beta$ le propone matrimonio a $C$.

\begin{figure}[H]\centering

\begin{tikzpicture}[ scale=0.8]
\tikzset{vertex/.style = {shape=circle,draw,minimum size=1.5em}}
\tikzset{edge/.style = {->,> = latex}}
\filldraw[color=blue!60, fill=blue!5, very thick](0,2.5) ellipse (.8 and 2);
\filldraw[color=red!60, fill=red!5, very thick](4,2.5) ellipse (.8 and 2);


% vertices
% 


\node[vertex] (a) at (0,4) {$\alpha$};
\node[vertex] (b) at (0,3) {$\beta$};
\node[vertex] (c) at (0,2) {$\gamma$};
\node[vertex] (d) at (0,1) {$\delta$};

\node[vertex] (e) at (4,4) {$A$};
\node[vertex] (f) at (4,3) {$B$};
\node[vertex] (g) at (4,2) {$C$};
\node [vertex] (h) at (4,1) {$D$};

\node (i) at (0,5) {Hombres};
\node (j) at (4,5) {Mujeres};

\path[-stealth] (a) edge (f);
\path[-stealth] (b) edge (g);
\path[-stealth] (c) edge (g);
\path[-stealth] (d) edge (e);



%\draw (0.2,8)--(3.8,8);



\end{tikzpicture}

\caption{Sexta iteración.}
\end{figure}

En séptima instancia, $C$ acepta la propuesta de $\beta$ y $\gamma$ le propone matrimonio a $A$.

\begin{figure}[H]\centering

\begin{tikzpicture}[ scale=0.8]
\tikzset{vertex/.style = {shape=circle,draw,minimum size=1.5em}}
\tikzset{edge/.style = {->,> = latex}}
\filldraw[color=blue!60, fill=blue!5, very thick](0,2.5) ellipse (.8 and 2);
\filldraw[color=red!60, fill=red!5, very thick](4,2.5) ellipse (.8 and 2);


% vertices
% 


\node[vertex] (a) at (0,4) {$\alpha$};
\node[vertex] (b) at (0,3) {$\beta$};
\node[vertex] (c) at (0,2) {$\gamma$};
\node[vertex] (d) at (0,1) {$\delta$};

\node[vertex] (e) at (4,4) {$A$};
\node[vertex] (f) at (4,3) {$B$};
\node[vertex] (g) at (4,2) {$C$};
\node [vertex] (h) at (4,1) {$D$};

\node (i) at (0,5) {Hombres};
\node (j) at (4,5) {Mujeres};

\path[-stealth] (a) edge (f);
\path[-stealth] (b) edge (g);
\path[-stealth] (c) edge (e);
\path[-stealth] (d) edge (e);



%\draw (0.2,8)--(3.8,8);



\end{tikzpicture}

\caption{Séptima iteración.}
\end{figure}


En octava instancia, $A$ acepta la propuesta de $\gamma$ y $\delta$ le propone matrimonio a $B$.

\begin{figure}[H]\centering

\begin{tikzpicture}[ scale=0.8]
\tikzset{vertex/.style = {shape=circle,draw,minimum size=1.5em}}
\tikzset{edge/.style = {->,> = latex}}
\filldraw[color=blue!60, fill=blue!5, very thick](0,2.5) ellipse (.8 and 2);
\filldraw[color=red!60, fill=red!5, very thick](4,2.5) ellipse (.8 and 2);


% vertices
% 


\node[vertex] (a) at (0,4) {$\alpha$};
\node[vertex] (b) at (0,3) {$\beta$};
\node[vertex] (c) at (0,2) {$\gamma$};
\node[vertex] (d) at (0,1) {$\delta$};

\node[vertex] (e) at (4,4) {$A$};
\node[vertex] (f) at (4,3) {$B$};
\node[vertex] (g) at (4,2) {$C$};
\node [vertex] (h) at (4,1) {$D$};

\node (i) at (0,5) {Hombres};
\node (j) at (4,5) {Mujeres};

\path[-stealth] (a) edge (f);
\path[-stealth] (b) edge (g);
\path[-stealth] (c) edge (e);
\path[-stealth] (d) edge (f);



%\draw (0.2,8)--(3.8,8);



\end{tikzpicture}

\caption{Octava iteración.}
\end{figure}

En novena instancia, $B$ acepta la propuesta de $\delta$ y $\alpha$ le propone matrimonio a $C$.

\begin{figure}[H]\centering

\begin{tikzpicture}[ scale=0.8]
\tikzset{vertex/.style = {shape=circle,draw,minimum size=1.5em}}
\tikzset{edge/.style = {->,> = latex}}
\filldraw[color=blue!60, fill=blue!5, very thick](0,2.5) ellipse (.8 and 2);
\filldraw[color=red!60, fill=red!5, very thick](4,2.5) ellipse (.8 and 2);


% vertices
% 


\node[vertex] (a) at (0,4) {$\alpha$};
\node[vertex] (b) at (0,3) {$\beta$};
\node[vertex] (c) at (0,2) {$\gamma$};
\node[vertex] (d) at (0,1) {$\delta$};

\node[vertex] (e) at (4,4) {$A$};
\node[vertex] (f) at (4,3) {$B$};
\node[vertex] (g) at (4,2) {$C$};
\node [vertex] (h) at (4,1) {$D$};

\node (i) at (0,5) {Hombres};
\node (j) at (4,5) {Mujeres};

\path[-stealth] (a) edge (g);
\path[-stealth] (b) edge (g);
\path[-stealth] (c) edge (e);
\path[-stealth] (d) edge (f);



%\draw (0.2,8)--(3.8,8);



\end{tikzpicture}

\caption{Novena iteración.}
\end{figure}


En decima instancia, $C$ acepta la propuesta de $\alpha$, $\beta$ le propone matrimonio a $D$ y $D$ acepta la propuesta de $\beta$. Después de 10 pasos en el algoritmo, todos los hombres ya están emparejados, es decir, el algoritmo convergió en un emparejamiento. Cabe mencionar que el número total de propuestas fue 13. 

\begin{figure}[H]\centering

\begin{tikzpicture}[ scale=0.8]
\tikzset{vertex/.style = {shape=circle,draw,minimum size=1.5em}}
\tikzset{edge/.style = {->,> = latex}}
\filldraw[color=blue!60, fill=blue!5, very thick](0,2.5) ellipse (.8 and 2);
\filldraw[color=red!60, fill=red!5, very thick](4,2.5) ellipse (.8 and 2);


% vertices
% 


\node[vertex] (a) at (0,4) {$\alpha$};
\node[vertex] (b) at (0,3) {$\beta$};
\node[vertex] (c) at (0,2) {$\gamma$};
\node[vertex] (d) at (0,1) {$\delta$};

\node[vertex] (e) at (4,4) {$A$};
\node[vertex] (f) at (4,3) {$B$};
\node[vertex] (g) at (4,2) {$C$};
\node [vertex] (h) at (4,1) {$D$};

\node (i) at (0,5) {Hombres};
\node (j) at (4,5) {Mujeres};

\path[-stealth] (a) edge (g);
\path[-stealth] (b) edge (h);
\path[-stealth] (c) edge (e);
\path[-stealth] (d) edge (f);



%\draw (0.2,8)--(3.8,8);



\end{tikzpicture}

\caption{Decima iteración.}
\end{figure}
\fin
\end{eje}

Lo imprtante de este algoritmo no es solo que produce un emparejamiento, adicional a esto como lo muestra el siguiente teorema, el emparejamiento resultante es estable. 
\begin{teo}[Teorema de Gale Shaley] \cite{GaleShapley} \\
\label{teorema de Gale Shapley}
El algoritmo de Gale Shapley termina en un emparejamiento estable.
\end{teo}
\begin{proof}
Supongamos que el emparejamiento producido por el algoritmo no es estable. Esto es, un hombre $\alpha$ y una mujer $A$ se prefieren entre ellos que a sus respectivas parejas. \\
Como $\alpha$ prefiere a $A$ más que a su esposa entonces, $\alpha$ le propuso matrimonio primero a $A$ que a su propia esposa. 
Además, como $A$ prefiere a $\alpha$ que a su esposo entonces, $A$ hubiera rechazado a su esposo y se hubiera quedado casada con $\alpha$ lo cual es una contradicción. \\
Por lo tanto, el algoritmo de Gale Shapley termina siempre en un emparejamiento estable. 
\end{proof}
Un resultado inmediato de esto es que siempre, sin importar como sean las preferencias, existe un emparejamiento estable. 
\begin{cor}
\label{corexiste}
Dada una matriz de preferencias arbitraria existe un emparejamiento estable. 
\end{cor}
\begin{proof}
Si aplicamos el algoritmo de Gale Shapley, sabemos por el teorema \ref{teorema de Gale Shapley} que este siempre acaba en un emparejamiento estable. Por lo tanto, siempre existe un emparejamiento estable.
\end{proof}

Una vez que sabemos que el algoritmo siempre converge, nos interesa conocer que tan rapido lo hace y si tiene algunas ventajas el emparejamiento resultante. El lema \ref{lema 1} nos ayuda a llegar a estos resultados y los corolarios \ref{cor1} y \ref{cor2} nos dan una idea de que tan bueno es el algoritmo. 
\begin{lem} 
\label{lema 1} \cite{Knuth} \\
Bajo el emparejamiento obtenido por Gale Shapley, solo un hombre puede terminar con la última mujer de su lista como pareja. 
\end{lem}

\begin{proof}

Supongamos que en el emparejamiento de Gale Shapley $m$ ($m\geq2$) hombres terminan con la última mujer de su lista como pareja, eso significa que cada uno de esos $m$ hombres invitó a salir a todas las mujeres. Entonces cada mujer fue invita a salir por lo menos $m$ veces, lo cual es una contradicción porque el algoritmo acaba cuando invitan a salir a la última mujer y a esta solo la invitan a salir una vez. Por lo tanto, solo un hombre puede terminar con la última mujer de su lista como pareja. 
\end{proof}

Dos consecuencias casi inmediatas de esto son que bajo el algoritmo la gran mayoria de los hombres no terminan con su ultima opción y que el algoritmo converge relativamente rapido. 
\begin{cor}
\label{cor1}
Si en un emparejamiento estable por lo menos dos hombres están emparejados con la última mujer de sus respectivas listas, entonces existe dos o más emparejamientos estables en el problema.
\end{cor}

\begin{proof}
Llamemos $m$ al emparejamiento estable donde por lo menos dos hombres están emparejados con la última mujer de sus respectivas listas y llamemos $m'$ al emparejamiento de Gale Shapley.
Por el teorema \ref{teorema de Gale Shapley} sabemos que el algoritmo de Gale Shapley siempre acaba en un emparejamiento estable y además en ese emparejamiento estable solo un hombre puede terminar con la última mujer de su lista como pareja por el lema \ref{lema 1}.
Por lo tanto $m$ y $m'$ son diferentes y como ambos son emparejamientos estables entonces el número de emparejamientos estables en el problema es mayor o igual a dos. 
\end{proof}

\begin{cor}
\label{cor2}
El número máximo de propuestas en el algoritmo es $n^2-n+1$
\end{cor}

\begin{proof}
Por el lema \ref{lema 1} sabemos que a lo más un hombre acaba emparejado con la última mujer de su lista, por lo tanto, el peor emparejamiento posible para el algoritmo es uno donde $n-1$ hombres terminan con la penúltima mujer de sus respectivas listas y un hombre termina con la última. Para llegar a esta situación de acuerdo con el algoritmo, los $n-1$ hombres deben de realizar $n-1$ propuestas cada uno y el otro hombre debe de realizar $n$ propuestas. Esto es $(n-1)(n-1)+n$ propuestas que es igual a $n^2-n+1$ propuestas.
\end{proof}

\begin{obs}
La complejidad del algoritmo de Gale Shapley es del orden de $n^2$.
\end{obs}

\begin{eje}
Si retomamos el ejemplo \ref{ejemploGS} podemos ver que para 4 personas hubo $4^{2}-4+1=13$ propuestas que es el máximo número posible de estas de acuerdo al corolario \ref{cor2}.
\fin
\end{eje}

A partir de esto ya podemos resolver el problema en el que la cantidad de hombres y de mujeres es distinta. El siguiente resultado muestra que sin importar el número de hombres o de mujeres, el algoritmo igual converge. 

\begin{teo} \cite{GaleShapley} \\
Dada una matriz de preferencias arbitraria con $n$ hombres y $m$ mujeres, el algortimo de Gale Shapley converge a un emparejamiento estable
\end{teo}
\begin{proof}
Supongamos que la cantidad de hombres es más chica que la cantidad de mujeres ($n<m$), en este caso el algoritmo acaba cuando $n$ de las $m$ mujeres reciben una propuesta. Si suponemos que la cantidad de hombres es mayor a la cantidad de mujeres ($n>m$), en este caso el algoritmo acaba despues de que $n-m$ hombres son rechazados por todas las mujeres y las propuestas de $m$ de los hombres son aceptadas. 

Además de forma analoga al teorema \ref{teorema de Gale Shapley} se puede ver que el emparejamiento producido es estable.
\end{proof}

El algoritmo de Gale Shapley no es el unico algoritmo que existe para producir un emparejamiento estable, y si se cambia el algoritmo el emparejamiento producido podría ser totalmente diferente al producido por el primero. De forma inmediata se puede crear un algoritmo igual al de Gale Shapley con la unica diferencia de que las mujeres le proponen a lo hombres, el siguiente ejemplo ilustra que el emparejamiento obtenido es distinto en ciertas situaciones y que por lo tanto los emparejamientos estables no son unicos. 

\begin{eje}
\label{ejeunico}
Retomando el ejemplo \ref{ejemplo matrimonio 1}. \\
Si aplicamos el algoritmo de Gale Shapley para los hombres obtenemos en primera estancia el siguiente emparejamiento estable. 
\begin{figure}[H]\centering

\begin{tikzpicture}[ scale=0.8]
\tikzset{vertex/.style = {shape=circle,draw,minimum size=1.5em}}
\tikzset{edge/.style = {->,> = latex}}
\filldraw[color=blue!60, fill=blue!5, very thick](0,3) ellipse (.8 and 1.5);
\filldraw[color=red!60, fill=red!5, very thick](4,3) ellipse (.8 and 1.5);


% vertices
% 


\node[vertex] (a) at (0,4) {$\alpha$};
\node[vertex] (b) at (0,3) {$\beta$};
\node[vertex] (c) at (0,2) {$\gamma$};


\node[vertex] (e) at (4,4) {$A$};
\node[vertex] (f) at (4,3) {$B$};
\node[vertex] (g) at (4,2) {$C$};


\node (i) at (0,5) {Hombres};
\node (j) at (4,5) {Mujeres};

\path[-stealth] (a) edge (e);
\path[-stealth] (b) edge (f);
\path[-stealth] (c) edge (g);



%\draw (0.2,8)--(3.8,8);



\end{tikzpicture}

\caption{Resultado del algoritmo si los hombres proponen.}
\end{figure}

Si aplicamos el algoritmo de Gale Shapley para las mujeres obtenemos en primera estancia el siguiente emparejamiento estable. 
\begin{figure}[H]\centering

\begin{tikzpicture}[ scale=0.8]
\tikzset{vertex/.style = {shape=circle,draw,minimum size=1.5em}}
\tikzset{edge/.style = {->,> = latex}}
\filldraw[color=blue!60, fill=blue!5, very thick](0,3) ellipse (.8 and 1.5);
\filldraw[color=red!60, fill=red!5, very thick](4,3) ellipse (.8 and 1.5);


% vertices
% 


\node[vertex] (a) at (0,4) {$\alpha$};
\node[vertex] (b) at (0,3) {$\beta$};
\node[vertex] (c) at (0,2) {$\gamma$};


\node[vertex] (e) at (4,4) {$A$};
\node[vertex] (f) at (4,3) {$B$};
\node[vertex] (g) at (4,2) {$C$};


\node (i) at (0,5) {Hombres};
\node (j) at (4,5) {Mujeres};

\path[-stealth] (g) edge (a);
\path[-stealth] (f) edge (b);
\path[-stealth] (e) edge (c);



%\draw (0.2,8)--(3.8,8);



\end{tikzpicture}

\caption{Resultado del algoritmo si las mujeres proponen.}
\end{figure}

Es facil ver que los dos emparejamientos estables son distintos.
\fin
\end{eje}

\begin{cor}
Dada una matriz de preferencias arbitraria, la cantidad de emparejamientos estables es mayor o igual a 1. 
\end{cor}
\begin{proof}
Con el corolario \ref{corexiste} podemos garantizar la existencia y con el ejemplo \ref{ejeunico} mostramos que no podemos garantizar la unicidad de este.
\end{proof}

Para continuar, podemos generalizar el problema del matrimonio estable a uno que permite a los hombres casarse con varias mujeres cambiando su cota superior de uno a cualquier otro número. Este es conocido como el problema de admisión a universidades. 

\section{El problema de admisión a universidades}
Supongamos que en una ciudad hay $n$ personas que desean entrar a $m$ universidades, los solicitantes tienen una lista de preferencias en donde reflejan a que universidades prefieren entrar y de forma analoga las universidades tienen una lista preferencias con la información de a quien prefieren admitir. Adicional a esto para cada universidad existe un numero máximo de alumnos que pueden admitir, esta restricción es natural porque la cantidad de personal y de espacio en las universidades es limitado. 
En terminos matemáticos, supongamos que tenemos $n$ solicitantes $$\alpha_1,\alpha_2,\ldots,\alpha_n$$ y $m$ universidades $$a_1, a_2,\ldots,a_m$$ con las restricciones
\begin{enumerate}
\item \begin{equation} \label{2r1}
x_{i,j}= 
\begin{cases}
1 & \qquad \text{si $i$ entra a estudiar a $j$.} \\
0 &\qquad\text{en otro caso.}\ \\ 
\end{cases} \end{equation}
\item \begin{equation} \label{2r2}
\sum_{j=1}^{m}x_{i,j} \leq1 \ \text{ para toda $i=1,2,\ldots,n$. }
\end{equation} Cada solicitante entra solo a una universidad. 
\item \begin{equation} \label{2r3}
\sum_{i=1}^{n} x_{i,j} \leq M_j\ \text{ para toda $j=1,2,\dots,m$.} 
\end{equation}
Cada universidad tiene un limite de alumnos que puede admitir. 
\end{enumerate}

Es claro que esto es una generalización del problema del matrimonio estable y un caso particular del problema de admisión a universidades con cotas inferiores y comunes. Podemos retomar las definiciones la matrz de preferencias, un emparejamiento estable y de un emparejamiento optimo como analogas a las definiciones \ref{matpref}, \ref{Estable} y \ref{optima}.

En el siguiente código exhibimos un análogo al algoritmo de Gale Shapley para este problema, el cual encuentra un emparejamiento estable y que además lo hace relativamente rapido. Para el caso general este algoritmo es conocido como el de ``'aceptación diferida''.
\pagebreak
\begin{lstlisting}[style=R, escapeinside={(*}{*)},caption={Algoritmo de Gale Shapley para admisión a universidades}, captionpos=b, label=c2]
Input : Una matriz de preferencias para (*$n$*) solicitantes y (*$m$*) universidades, un vector (*$M$*) en donde la entrada (*$j$*) representa la cota superior de la universidad (*$j$*).
Output: Un emparejamiento. 
Cada solicitante aplica a la primera universidad de su lista. Cada universidad que recibe más de solicitudes a su cota superior acepta las que solicitudes que estan más arriba en su lista y rechaza al resto. 
repeat{ #hasta que cada uno de los solicitantes sea admitido por alguna universidad o rechazado por todas.
	Los solicitantes no emparejados solicitan entrar a la siguiente universidad en su lista. 
	Cada universidad que recibe alguna solicitud acepta hasta (*$M_j$*) solicitantes de los primeros de su lista entre los que aplicaron a ella y sus admitidos actuales.
	Las universidades rechazan a los alumnos no aceptados.
}
\end{lstlisting}

Este algoritmo al igual que su análogo en el matrimonio estable encuentra siempre un emparejamiento estable.
\begin{cor}
Dada una matriz de preferencias arbitraria y un vector $M$ de cotas superiores arbitrario, el algoritmo de Gale Shapley siempre encuentra un emparejamiento estable.
\end{cor}
\begin{proof}
Al igual que en la demostración del teorema \ref{teorema de Gale Shapley}, supongamos que el emparejamiento producido por el algoritmo no es estable. 
Esto es, un solicitante $\alpha$ que esta admitido en una universidad $B$ (o en ninguna) prefiere estar en una universidad $A$ y simultáneamente una universidad $A$ tiene admitido a un alumno $\beta$ y preferiría tener a $\alpha$ que a $\beta$ como estudiante. \\
Como $\alpha$ prefiere a $A$ más que a su universidad entonces, $\alpha$ solicito entrar primero a $A$ que a su propia universidad. 
Además, como $A$ prefiere a $\alpha$ que a $\beta$ entonces de acuerdo con el algoritmo, $A$ hubiera rechazado a $\beta$ y se hubiera quedado con $\alpha$ como alumno lo cual es una contradicción (el argumento para cuando $\alpha$ no fue aceptado en ninguna universidad es el mismo). \\
Por lo tanto, el algoritmo de Gale Shapley termina siempre en un emparejamiento estable. 
\end{proof}

Este algoritmo además de ser producir un emparejamiento estable también cada solicitante esta mejor o igual en este emparejamiento que en cualquier otro emparejamiento estable. Para hacer la demostración primero introducimos una definición y un lema. 

\begin{dfn}{\cite{GaleShapley}}
\label{Posible}
Decimos que una universidad es \textbf{posible} para un aplicante si existe una asignación estable en la que esta persona asiste a esa universidad.
\end{dfn}

\begin{lem} 
\label{lema optimo} 
\cite{GaleShapley} \\
Supongamos que en un paso arbitrario del algoritmo ningun estudiante a sido rechazado por una universidad posible para el, además supongamos que una universidad $A$ despues de llenarse recibiendo a los estudiantes $\beta_1,\beta_2,\dots,\beta_q$ rechaza a $\alpha$, entonces $A$ no es posible para $\alpha$.
\end{lem}
\begin{proof}
Sabemos por hipotesis que para toda $i=1,\dots,q$, $\beta_i$ prefiere a $A$ que a todas las universidades que no lo han rechazado y además que cualquier universidad que lo rechazo previamente no es posible para el. Supongamos que existe un emparejamiento estable en el que $\alpha$ asiste a $A$, entonces  alguna $\beta_i$ no asiste a $A$ porque $\alpha$ tomo su lugar. Este emparejamiento es inestable porque $\beta_i$ prefiere a $A$ que a su asignación actual porque $A$ es su mejor asignación posible y $A$ prefiere tener a $\beta_i$ que a $\alpha$, lo cual es claramente una contradicción y por lo tanto $A$ no es posible para $\alpha$.
\end{proof}

\begin{teo}
\label{optimo}
\cite{GaleShapley}  \\
Dada una matriz de preferencias arbitraria y un vector de cotas superiores arbitrario, el emparejamiento producido por el algoritmo es óptimo para los solicitantes.
\end{teo}
\begin{proof}
La prueba es por inducción. Primero que nada sabemos que si en el primer paso del algoritmo una universidad rechaza a un alumno es porque esta universidad no es posible para el aplicante, si suponomes que esta universidad es posible para el llegamos rapidamente a una contradicción porque esto quisiera decir que la universidad rechazo a un mejor estudiante que la tenia como primer opción para meterlo a el y por lo tanto el emparejamiento no seria estable. Luego por el lema \ref{lema optimo} sabemos que en los siguientes pasos del algoritmo ningun estudiante es rechazado por una universidad posible para el y por lo tanto el emparejamiento obtenido es optimo. 
\end{proof}

El algoritmo mencionado anteriormente no es unico, de cambiarlo el emparejamiento obtenido podría cambiar significativamente. En primera instancia el algoritmo puede ser modificado para que sea optimo para las universidades en lugar de para los solicitantes. El siguiente algoritmo produce este resultado y las demostraciones se hacen de forma analoga a las mencionadas anteriormente, como propiedades importantes vale la pena mencionar que produce un emparejamiento estable y que cada universidad esta mejor en este emparejamiento que en cualquier otro emparejamiento. 

\begin{lstlisting}[style=R, escapeinside={(*}{*)},caption={Algoritmo de Gale Shapley para admisión a universidades modificiado}, captionpos=b, label=c3]
Input : Una matriz de preferencias para (*$n$*) solicitantes y (*$m$*) universidades, un vector (*$M$*) en donde la entrada (*$j$*) representa la cota superior de la universidad (*$j$*).
Output: Un emparejamiento. 
Cada universidad invita a los primeros (*$j_i$*) solicitantes en su lista a estudiar ahí. Cada alumno que recibe más de una solicitud acepta las más alta en su lista y rechaza el resto. 
repeat{ #hasta que cada universidad este llene o cuando todos los estudiantes sean admitidos en una universidad.
	Las universidades que no alcanzaron su cota superior invitan a los siguientes alumnos de su lista de acuerdo a cuantos lugares tienen.
	Cada solicitante que recibe una solicitud acepta la más alta en su alta entre su acutal universidad y las que lo invitaron. 
	El alumno rechaza el resto de las invitaciones.
}
\end{lstlisting}

\subsection{El teorema de los hospitales rurales}

Existe una extensión de este problema llamada el problema de hospitales y residentes medicos que es basicamente lo mismo a este problema con la excepción de que en vez de tener universidades tenemos hospitales y en vez de que existan estudiantes queriendo estudiar en esta tenemos estudiantes medicos que quieren hacer su residencia en este. Una pregunta interesante que se hizo fue ¿qué pasa con los hospitales no populares? Algunos hospitales por su naturaleza no llamaban la atención de los estudiantes de medicina, un ejemplo de esto eran los que estaban en una zona rural, los residentes preferian ir a hospitales en zonas urbanas más pobladas porque contaban con fama de ser mejores hospitales y además tenian más pacientes. La pregunta original puede ser modificada en ¿existe una modificación al algoritmo que produzca un emparejamiento estable y además asigne más residentes en los hospitales poco populares que el resultado de el algortimo Gale Shapley? A raíz de esto David Gale y Marilda Sotomayor en 1985 demostraron que la cantidad de gente que asiste a universidad es la misma en cualquier emparejamiento estable. Este resultado se puede poner en terminos de nuestro problema original, para llegar a este es necesario primero demostrar un lema.

\begin{lem} 
\label{lema rural} 
\cite{Verde} \\
Dada una matriz de preferencias arbitraria y un vector de cotas superiores arbitario, sea $M$ el emparejamiento obtenido por el algoritmo de Gale Shapley para admisión a universidades y sea $M'$ un emparejamiento estable arbitrario. Si una universidad $A$ no se llena en $M'$ entonces todo solicitante que entro a $A$ en $M$ tambien fue asignado a $A$ en $M'$. 
\end{lem}
\begin{proof}
Supongamos que $\alpha$ fue admitido en $A$ en la asignación $M$ pero no en $M'$, por hipotesis $A$ no esta llena en $M'$ entonces tenemos una contradicción porque significaria que $\alpha$  en $M'$ esta en una mejor universidad posible que $A$, lo cual no es posible porque la asignación en $M$ por el teorema \ref{optimo} es optima. Por lo tanto  Si una universidad $A$ no se llena en $M'$ entonces todo solicitante que entro a $A$ en $M$ tambien fue asignado a $A$ en $M'$. 
\end{proof}


\begin{teo}{Teorema de los hospitales rurales \\ }
\label{rural}
\cite{GaleSotomayor}\\
Dada una matriz de preferencias arbitraria y un vector de cotas superiores arbitario,:
\begin{enumerate}
\item Cada universidad recibe el mismo número de solitantes en cada asignación estable.
\item Exactamente el mismo número de solicitantes quedan como no asignados en cada emparejamiento estable. 
\item Si una universidad no se llena en un emparejamiento estable recibe exactamente el mismo número de solicitantes en cualquier asignación estable. 
\end{enumerate}
\end{teo}
\begin{proof}
Sea M el emparejamiento obtenido por Gale Shapley y sea $M'$ un emparejamiento estable arbitario. Prime notamos que si una persona no fue asignada a ningun universidad en $M$ entonces tampoco en $M'$, esto es por el lema \ref{lema optimo} porque ninguna universidad es posible para el. Por lo tanto el número de personas asignadas en $M'$ no puede exceder el de $M$. \\
Ahora por el lema \ref{lema rural}, si una universidad se llena en $M$ entonces tambien esta llena en $M'$.  Simultaneamente por el lema \ref{lema rural}, si a una universidad le sobran lugares en $M'$ entonces en $M$ tiene minimo la misma cantidad de lugares. Por lo tanto el número de personas asignadas en $M$ no puede exceder el de $M'$. Luego, podemos concluir que para toda universidad  tiene el mismo número de estudiantes asignados en $M$ y en $M'$ que es lo mismo a cada universidad recibe el mismo número de solitantes en cada asignación estable. \\
Una conclunsión imediata de arriba y de que ningun estudiante que es no asignado en $M$ puede ser admitido en $M'$ es que exactamente el mismo número de solicitantes quedan como no asignados en cada emparejamiento estable. \\
Una segunda conclusión inmediata de arriba y del lema \ref{lema rural} es que si una universidad no se llena en un emparejamiento estable recibe exactamente el mismo número de solicitantes en cualquier asignación estable.\\
Por lo tanto concluimos que se cumplen 1, 2 y 3.
\end{proof}

Ahora podemos retomar el problema original y ver si podemos extender o usar los resultados obtenidos. 

\section{El problema de admisión a universidades con cotas inferiores y comunes}

Una pregunta inicial que surge es que si todos los resultados obtenidos para el problema sin cotas inferiores o comunes se siguen cumpliendo ¿el emparejamiento estable siempre existe? ¿De existir este es optimo? ¿Se obtiene con facilidad? Para comenzar es posible mostrar que el emparejamiento estable no siempre existe.

\begin{eje}
\cite{Todo}\\
Supongamos que tenemos dos universidades $A$ y $B$, $A$ solo puede admitir a dos alumnos y necesita admitir a dos alumnos para entrar y $B$ solo puede admitir a un alumno y necesita a un alumno para abrir. Además contamos con dos solicitantes $\alpha$ y $\beta$, la matriz de preferencias del problema esta dada por 
$$\begin{pmatrix}
& A & B  \\
\alpha & 1,1 & 2,1  \\
\beta & 2,2 & 1,2 \\ 
\end{pmatrix}.$$
Lo primero que notamos es que $A$ o $B$ tienen que estar cerradas porque la suma de las cotas inferiores es tres y solo contamos con dos estudiantes. Si $A$ cierra entonces tenemos dos emparejamientos posibles, en el primero $\alpha$ es admitido en $B$ y no es estable porque $\alpha$ prefiere estar en $A$, $\beta$ prefiere estar en $A$ a no estar en ningun lado y $A$ prefiere tener a los dos estudiantes que cerrar . 
\begin{figure}[H]\centering

\begin{tikzpicture}[ scale=0.8]
\tikzset{vertex/.style = {shape=circle,draw,minimum size=1.5em}}
\tikzset{edge/.style = {->,> = latex}}
\filldraw[color=blue!60, fill=blue!5, very thick](0,3) ellipse (.8 and 1.5);
\filldraw[color=green!60, fill=green!5, very thick](4,3) ellipse (.8 and 1.5);


% vertices
% 


\node[vertex] (a) at (0,4) {$\alpha$};
\node[vertex] (b) at (0,2) {$\beta$};


\node[vertex] (e) at (4,4) {$A$};
\node[vertex] (f) at (4,2) {$B$};


\node (i) at (0,5) {Solicitantes};
\node (j) at (4,5) {Universidades};

\path[-stealth] (a) edge (f);
%\path[-stealth] (b) edge (g);

\draw (4,4) node[cross,red] {};


%\draw (0.2,8)--(3.8,8);



\end{tikzpicture}

\caption{Primer caso.}
\end{figure}
En el segundo caso $\beta$ es admitido en $B$ y esta emparejamiento no es estable porque $B$ prefiere tener a $\alpha$ que a $\beta$ y $\alpha$ prefiere estudiar que no estudiar. 
\begin{figure}[H]\centering

\begin{tikzpicture}[ scale=0.8]
\tikzset{vertex/.style = {shape=circle,draw,minimum size=1.5em}}
\tikzset{edge/.style = {->,> = latex}}
\filldraw[color=blue!60, fill=blue!5, very thick](0,3) ellipse (.8 and 1.5);
\filldraw[color=green!60, fill=green!5, very thick](4,3) ellipse (.8 and 1.5);


% vertices
% 


\node[vertex] (a) at (0,4) {$\alpha$};
\node[vertex] (b) at (0,2) {$\beta$};


\node[vertex] (e) at (4,4) {$A$};
\node[vertex] (f) at (4,2) {$B$};


\node (i) at (0,5) {Solicitantes};
\node (j) at (4,5) {Universidades};

\path[-stealth] (b) edge (f);
%\path[-stealth] (b) edge (g);

\draw (4,4) node[cross,red] {};


%\draw (0.2,8)--(3.8,8);



\end{tikzpicture}

\caption{Segundo caso.}
\end{figure}


Si $B$ cierra entonces solo tenemos un emparejamiento posible en el que $\alpha$ y $\beta$ son admitidos en $A$, este no es estable porque $B$ prefiere tener a $\beta$ que cerrar y $\beta$ prefiere estar en $A$ que en $B$.

\begin{figure}[H]\centering

\begin{tikzpicture}[ scale=0.8]
\tikzset{vertex/.style = {shape=circle,draw,minimum size=1.5em}}
\tikzset{edge/.style = {->,> = latex}}
\filldraw[color=blue!60, fill=blue!5, very thick](0,3) ellipse (.8 and 1.5);
\filldraw[color=green!60, fill=green!5, very thick](4,3) ellipse (.8 and 1.5);


% vertices
% 


\node[vertex] (a) at (0,4) {$\alpha$};
\node[vertex] (b) at (0,2) {$\beta$};


\node[vertex] (e) at (4,4) {$A$};
\node[vertex] (f) at (4,2) {$B$};


\node (i) at (0,5) {Solicitantes};
\node (j) at (4,5) {Universidades};

\path[-stealth] (a) edge (e);
\path[-stealth] (b) edge (e);

\draw (4,2) node[cross=10pt,red] {};


%\draw (0.2,8)--(3.8,8);



\end{tikzpicture}

\caption{Tercer caso.}
\end{figure}

\fin
\end{eje}

A partir de este ejemplo queda claro que si extendemos el problemo agregando cotas inferiores o comunes perdemos la garantia de existencia de un emparejamiento estable en nuestro problema. Es posible además consturir un algoritmo a partir del algoritmo de Gale Shapley para admisión a universidades que converja en caso de que exista a un emparejamiento estable o que saque como resultado si no existe un emparejamiento estable para este problema, el algoritmo esta dado por el siguiente codigo.
 
 \begin{lstlisting}[style=R, escapeinside={(*}{*)},caption={Algoritmo de Gale Shapley para admisión a universidades modificiado}, captionpos=b, label=c3]
Input : Una matriz de preferencias para (*$n$*) solicitantes y (*$m$*) universidades, un vector (*$M$*) en donde la entrada (*$j$*) representa la cota superior de la universidad (*$j$*) , un vector $N$ en donde la entrada (*$j$*) representa la cota superior de la universidad (*$j$*) y una matriz $$.
Output: Un emparejamiento. 
Cada universidad invita a los primeros (*$j_i$*) solicitantes en su lista a estudiar ahí. Cada alumno que recibe más de una solicitud acepta las más alta en su lista y rechaza el resto. 
repeat{ #hasta que cada universidad este llene o cuando todos los estudiantes sean admitidos en una universidad.
	Las universidades que no alcanzaron su cota superior invitan a los siguientes alumnos de su lista de acuerdo a cuantos lugares tienen.
	Cada solicitante que recibe una solicitud acepta la más alta en su alta entre su acutal universidad y las que lo invitaron. 
	El alumno rechaza el resto de las invitaciones.
}
\end{lstlisting}



%Falta arreglar definición de estabilidad

\thispagestyle{empty}
\chapter{Problemas NP-Completos}

\section{Clases de complejidad}
Algo que nos interesa mucho cuando atacamos un problema de computación es que tan dificíl es de resolver y si existe una mejor manera de hacerlo. La complejidad computacional nos permite estudiar esto. Para encarrilarnos hay que empezar con algunas definiciones y ejemplos.
\begin{dfn}
Se dice que un problema es un problema de decisón si admite dos posibles respuestas ``Si'' o ``No''. Resolverlo es encontrar todas las soluciones en las que la respuesta es ``Si''.
\end{dfn}
Existen dos clases de problemas de decisión que nos interesan, los que son faciles de verificar (NP) y los que son faciles de resolver (P)-
\begin{dfn}
Un problema esta en la clase NP si existe un algoritmo polinomial para verificar que la respuesta de una solución dada es ``Si''.
\end{dfn}

\begin{eje}[Problema del clima]
Supongamos que Mariana le dice a Adrián que esta lloviendo, si Adrián decide verificar si esto es cierto solo tiene que realizar tres pasos:
\begin{enumerate}
\item Caminar a la ventana.
\item Abrir la ventana.
\item Ver el cielo.
\end{enumerate}
Por lo tanto a Adrián le toma un tiempo constante verificar si esta lloviendo lo que implica que este problem esta en NP.
\end{eje}
\begin{dfn}
Se dice que un problema de decisión esta en P si existe un algoritmo polinomial que lo resuelve.
\end{dfn}
\begin{obs}[$P  \subseteq NP$]
Si un problema de decisión pertenece a P entonces pertenece a NP.
\end{obs}
%\begin{proof}
%Supongamos que $A$ es un problema de decision en P y $b$ es una solución dada de $A$, podemos verificar si la respuesta de la solución es ``Si'' resolviendo el problema en tiempo polinomial.
%\end{proof}
\begin{obs}
\label{pp eq}
Si A es un problema de decisión en P y existe un algoritmo polinomial que reduce resolver el problema de decisión B a resolver el problema A entonces B esta en P.
\end{obs}

Vale la pena mencionar algunos ejemplos de problemas que estan en P:
\begin{enumerate}
\item El problema de admisión a universidades.
\item Verificar si un numero es primo.
\item Encontrar la ruta más corta entre dos vertices en una gráfica.
\item El problema de la ruta critica. 
\end{enumerate}

Una conjetura famosa que vale la pena mencionar es ¿si $P=NP$?, es decir si todo problema en NP tambien se enecuentra en P. La persona que lo demuestra además de ganar fama y trabajo va a ser acredor a un permio de un millon de dolares. 

\section{El $3-SAT$}
Un problema de gran importancia para entender las clases de complejidad es el $3-SAT$, la raíz de todo el estudio necesario para estudiar el problema ¿$P=NP?$ se remonta a enteder este problema. 
\begin{dfn} [conjunctive normal form.]
\end{dfn}


\section{Problemas NP-Completos}
\begin{flushright}
\textit{``Chaos is a ladder'}
\end{flushright}
\begin{dfn}
Qué es un problema NP-completo. 
\end{dfn}
\begin{teo}[Teorema de Cook-Levin]
El 3-SAT es NP-Completo.
\end{teo}
Donde esta la demostración. 

\thispagestyle{empty}
\chapter{Conclusiones}

En esta tesis se construyó el camino para explicar los problemas de admisión a universidades con cotas inferiores o con cotas comunes. Estos a pesar de su facilidad al momento de explicarlos tienen la propiedad de ser NP-Completos y por lo tanto no son fáciles de resolver. Al mismo tiempo se vio que si se hace una reducción al problema de cotas comunes, esté es fácil de resolver y simultáneamente notamos que cuenta con una estructura algebraica muy poderosa. 

Al principio de la tesis se habló de como estos problemas se pueden plantear usando restricciones enteras, esto fue con el fin de plantear todo en el mundo de investigación de operaciones en el que estamos todos acostumbrados. Además, se aprovechó para dar una muy breve introducción a las clases de complejidad por si el lector no se encontraba familiarizado con el tema. 

Posteriormente, se dio una introducción a dos problemas realmente simples: el del matrimonio estable y el de admisión de universidades. Estos a pesar de su simplicidad son el primer paso para entender la situación completa de la tesis y nos dieron herramientas para estudiar el resto de los resultados. 

Aprovechando los resultados anteriores se plantearon tres problemas y se demostró que dos de ellos son NP-Completos (el tercero también cumple esta propiedad, pero se decidió que la demostración no iba con los objetivos de la tesis). Estos problemas resultan interesantes porque con solo incluir una restricción a los problemas originales estos pasan de ser muy fáciles de resolver a muy difíciles. 

Al final se vio que uno de los tres problemas de los planteados (el de cotas comunes) si le agregamos la propiedad de que estas son anidadas, el problema pasa de regreso a ser muy simple de resolver. Además, usando una estructura algebraica llamada matroides tuvimos la oportunidad de reevaluar muchos de los resultados de la tesis en un contexto diferente. 

\thispagestyle{empty}


%----------------------------------------------------------------------------------------
%	APÉNDICES
%----------------------------------------------------------------------------------------

\addtocontents{toc}{\vspace{2em}} % Agrega espacios en la toc

\appendix % Los siguientes capítulos son apéndices

%  Incluye los apéndices en el folder de apéndices

\chapter{Álgebra Lineal}

Cosas.
\thispagestyle{empty}
%\chapter{Teoría de Gráficas}

\begin{dfn}  
Una \textbf{gráfica} es una pareja $G = (V, E)$ donde $V$ es un conjunto finito no vacío
cuyos elementos se llaman \textbf{vértices} y $E$ es un conjunto cuyos elementos son subconjuntos de
cardinalidad 2 de V y estos son llamados \textbf{aristas}.
\end{dfn}

\begin{dfn}  
El \textbf{grado} de un vértice $v$ es el número de aristas que inciden en $v$, se denota $d(v)$.
\end{dfn}

\begin{dfn} 
Se dice que una gráfica es \textbf{bipartita}, si su conjunto de vértices puede ser partido en dos subconjuntos ajenos $X$ y $Y$, de tal forma que cada arista tiene un extremo en $X$ y otro en $Y$. Esta partición tiene el nombre de $\textbf{bipartición}$ de la gráfica.
\end{dfn}


%\include{Apendices/AppendixC}

\addtocontents{toc}{\vspace{2em}} % Agrega espacio en la toc


%----------------------------------------------------------------------------------------
%	BIBLIOGRAFÍA
%----------------------------------------------------------------------------------------

\backmatter
\nocite{*}
\bibliographystyle{apacite}
\bibliography{tesis_itam_bibliografia} %Aquí ponen el nombre del archivo .bib



\end{document}